% =====================================================================
%  LAB 5 -- Smoothing, Laplace, Unsharp Masking, GPU Challenge
%  Headings copied word for word from DSP2_Lab5_Exercise-2.pdf
%  NOTE: the Lab 5 sheet has no single topic line under "Lab Exercise",
%  so the chapter title below combines the four task titles -- rename it
%  if you prefer something else.
% =====================================================================
\clearpage
\section{Lab exercise 5: Smoothing and Noise Reduction, Edge Enhancement, Unsharp Masking, GPU Challenge}

TODO: chapter introduction -- theory about smoothing filters, Laplace
operator, unsharp masking and GPU computing.

\subsection{Laboratory 5 -- Smoothing and Noise Reduction}

% description given on the exercise sheet:
Smoothing filters are used to reduce noise, but they often result in blurring

\subsubsection{(a) Manual Implementation: Import the image 'cameraman.tif'. Implement a 3\hx3 Box Filter using only for loops. Each pixel should be replaced by the average of its 3\hx3 neighborhood.}

TODO

\subsubsection{(b) Comparison: Apply a 3\hx3 Gaussian Filter using the kernel G provided in the sources:}

% kernel given on the exercise sheet:
\[
  G = \frac{1}{16}
  \begin{pmatrix}
    1 & 2 & 1\\
    2 & 4 & 2\\
    1 & 2 & 1
  \end{pmatrix}
\]
Explain why the Gaussian filter results in less contour blurring than the Box filter.

TODO

\subsubsection{(c) Boundary Handling: Briefly describe how your code handles the edges of the image (e.g., zero-padding or ignoring borders).}

TODO

\subsection{Laboratory 5 -- Edge Enhancement with the Laplace Operator}

% description given on the exercise sheet:
The Laplace operator is used to highlight intensity changes by calculating the
second derivative.

\subsubsection{(a) Implementation: Apply the standard 3\hx3 Laplace kernel \texorpdfstring{$H_L$}{HL} to the Y (Luminance) component of the image 'peppers.png':}

% kernel given on the exercise sheet:
\[
  H_L =
  \begin{pmatrix}
    0 & 1 & 0\\
    1 & -4 & 1\\
    0 & 1 & 0
  \end{pmatrix}
\]
Use manual convolution (nested loops).

TODO

\subsubsection{(b) Visualization: The resulting image contains negative values. Explain how you visualize this ``Laplace image''correctly (e.g., using absolute values or shifting the zero-point).}

TODO

\subsection{Laboratory 5 -- Unsharp Masking}

% description given on the exercise sheet:
Unsharp masking enhances edges by subtracting a blurred version of the image
from the original

\subsubsection{(a) Calculation: Using your results from Task 1 and 2, implement a simple sharpening step: \texorpdfstring{$I_{sharp} = I_{original} - \mathrm{Laplace}(I_{original})$}{Isharp = Ioriginal - Laplace(Ioriginal)}.}

TODO

\subsubsection{(b) Effect: Describe the visual difference between the original image and the sharpened result.}

TODO

\subsection{Laboratory 5 -- GPU Challenge}

\subsubsection{(a) Benchmark a 5\hx5 convolution on a large image (e.g., 4000\hx4000) generated by the fractal.m. Compare the execution time of your nested for loops on the CPU versus the parallel execution on a GPU using gpuArray.}

% hints given on the exercise sheet:
\textbf{Hints for MATLAB GPU Implementation} To successfully complete the GPU
challenge, consider the following technical steps in MATLAB:
\begin{itemize}
  \item \textbf{Data Transfer:} Use the function \texttt{gpuArray(A)} to
        transfer a matrix from the CPU memory (RAM) to the GPU memory (VRAM).
  \item \textbf{Synchronization:} GPU operations are often asynchronous. To get
        an accurate measurement using \texttt{tic} and \texttt{toc}, you must
        call \texttt{wait(gpuDevice)} before stopping the timer to ensure all
        parallel threads have finished their execution.
  \item \textbf{Retrieval:} After the computation is complete, use the
        \texttt{gather(G)} function to move the result back to the CPU memory
        so it can be processed by standard functions like \texttt{imshow} or
        \texttt{imwrite}.
  \item \textbf{Compatibility:} You can check if a compatible hardware device
        is available by using the \texttt{gpuDeviceCount} function.
\end{itemize}

TODO
